% Instructor-authored original practice bank. This is an includable LaTeX fragment.
% Each source-map comment records only the disclosed form location and an abstract skill.

\section*{2025 Reading Comprehension: Original Analogues}

\begin{center}
\fbox{\parbox{0.92\linewidth}{\textbf{Instructor-authored original material.}
The names LSAT and LSAC are used only to identify general skill inspiration.
Every passage, question, and answer choice below is newly written.
This material is not written, approved, or endorsed by LSAC.}}
\end{center}

\noindent This is a public practice bank, not a live graded form. Each packet is designed for approximately 30--50 minutes of reading, individual analysis, and group discussion.

\subsection*{Packet 1: Mapping Sound}

City maps ordinarily privilege what can be seen: streets, buildings, property lines, and transit routes. When planner Imani Vale began documenting neighborhood sound in the city of Bellwether, some colleagues treated her project as decorative. At best, they thought, a sound layer might help tourists find outdoor concerts. Vale instead argued that sound maps could reveal barriers that conventional maps obscure. A crossing signal that lacks an audible cue may be unusable for a blind pedestrian; a supposedly quiet courtyard may sit beneath ventilation equipment whose low hum is exhausting to nearby residents; a public meeting room may transmit speech poorly even when its seating plan appears accessible.

Vale's first maps combined short recordings with measurements taken at different times of day. Crucially, she also invited residents to annotate the maps. A measured volume alone could not distinguish a welcome playground from an intrusive loading dock, and it could not show which sounds carried practical information. Residents' descriptions supplied this context. The annotations sometimes disagreed: one person's lively market was another person's obstacle to conversation. Vale did not regard that disagreement as a defect. For her, it showed that acoustic planning involves judgments about use, not merely readings from instruments.

The transportation department initially resisted Vale's recommendation that sound be considered during project review. Officials worried that because reactions differ, acoustic evidence could never support consistent decisions. Vale proposed a narrower rule. Projects would not be required to maximize quiet. Instead, reviewers would ask whether a project removed an important audible signal, imposed a persistent sound on people who could not readily avoid it, or excluded a group from participating in a public space. These questions, she maintained, could be answered with evidence even when no universal ideal soundscape existed.

Bellwether eventually added an acoustic section to its review forms. The change did not settle every dispute, and Vale never claimed it would. Its achievement was institutional: sound ceased to be an afterthought raised only after construction. Designers began discussing it while alternatives were still available. The maps mattered, then, not because they converted experience into a perfectly objective score, but because they made an overlooked dimension of access available for public reasoning.

\begin{enumerate}

% Source map: 2025 RC, Question 01 — main point
\item Which statement best expresses the main point of the passage?
\begin{enumerate}
\item[(A)] Sound maps can replace conventional city maps whenever accessibility is the primary concern.
\item[(B)] Although acoustic experience cannot be reduced to one objective score, documenting it can make overlooked access issues part of timely public decisions.
\item[(C)] Residents generally agree about which urban sounds are desirable.
\item[(D)] Bellwether adopted sound review chiefly to attract tourists to concerts.
\item[(E)] Instrument measurements should never be used in acoustic planning.
\end{enumerate}

% Source map: 2025 RC, Question 02 — stated detail EXCEPT
\item According to the passage, resident annotations helped do each of the following \emph{except}
\begin{enumerate}
\item[(A)] distinguish sounds experienced as welcome from sounds experienced as intrusive.
\item[(B)] identify sounds that carried practical information.
\item[(C)] make disagreement among residents visible.
\item[(D)] determine the precise calibration settings of the recording instruments.
\item[(E)] supply context that measured volume alone could not provide.
\end{enumerate}

% Source map: 2025 RC, Question 03 — viewpoint summary
\item Which statement best summarizes Vale's position on differing reactions to the same sound?
\begin{enumerate}
\item[(A)] Differing reactions prove that acoustic evidence should be excluded from review.
\item[(B)] Differing reactions matter, but they do not prevent reviewers from answering narrower evidence-based questions about access.
\item[(C)] Differing reactions disappear when volume is measured accurately.
\item[(D)] Only the reaction of the quietest resident should guide a decision.
\item[(E)] Every reaction should be converted into the same numerical score.
\end{enumerate}

% Source map: 2025 RC, Question 04 — function of a qualification
\item The statement that Vale ``never claimed'' the policy would settle every dispute primarily serves to
\begin{enumerate}
\item[(A)] concede a limit while preserving the passage's claim about the policy's institutional value.
\item[(B)] show that Vale considered the policy a failure.
\item[(C)] prove that officials were correct to reject all acoustic evidence.
\item[(D)] introduce a dispute about the accuracy of recording devices.
\item[(E)] suggest that the policy was adopted without Vale's knowledge.
\end{enumerate}

% Source map: 2025 RC, Question 05 — viewpoint application
\item Vale would most likely approve of which project-review practice?
\begin{enumerate}
\item[(A)] Assigning every site a single acceptable decibel level regardless of its use
\item[(B)] Ignoring resident reports whenever two reports conflict
\item[(C)] Waiting until a building opens before considering whether announcements can be heard
\item[(D)] Testing whether a proposed plaza preserves audible navigation cues and permits conversation before finalizing its design
\item[(E)] Prohibiting every sound that some resident dislikes
\end{enumerate}

% Source map: 2025 RC, Question 06 — shared support
\item The passage provides the strongest support for which statement about both measurements and resident annotations?
\begin{enumerate}
\item[(A)] Each is treated as useful evidence, but neither alone fully captures how sound affects access and use.
\item[(B)] Each produces a perfectly objective ranking of proposed projects.
\item[(C)] Both were rejected by the transportation department.
\item[(D)] Both are valuable only after construction is complete.
\item[(E)] Each eliminates the need for public judgment.
\end{enumerate}

\end{enumerate}

\subsection*{Packet 2: Two Views of Urban Trees}

\subsubsection*{Passage A}

Urban tree programs often celebrate a few immense specimens: the elm beside city hall, the oak pictured on neighborhood signs, the plane trees shading an old boulevard. This emphasis is understandable but risky. A canopy dominated by a small number of closely related trees can be devastated by a single pest. Even when a famous tree is healthy, maintaining it may consume funds that could establish dozens of younger trees in hotter, underserved blocks.

Cities should therefore evaluate tree programs at the scale of the whole canopy. Diversity of species and age makes the system less vulnerable, while broad geographic distribution brings shade and cleaner air to more residents. This does not mean that an old tree has no value. It means only that sentimental attachment should not automatically outweigh the resilience and fairness of the larger program. The most responsible memorial to a beloved tree may sometimes be a diverse grove planted where canopy is scarce.

\subsubsection*{Passage B}

Calls to manage the ``whole canopy'' can conceal what is lost when veteran trees are treated as interchangeable units. An old tree is not simply a large version of a young one. Hollows that shelter animals form slowly; extensive roots moderate water across a wide area; generations of residents attach stories to a familiar landmark. Planting twenty saplings does not instantly reproduce these ecological and cultural functions.

This is not an argument for uniformity. New planting should absolutely diversify urban forests, especially where shade is scarce. But diversity planning and preservation need not be rivals. A city can protect veteran trees while adding other species around them and across neglected neighborhoods. Removal should be a last resort supported by a specific hazard or an unavoidable conflict, not the default consequence of a spreadsheet that counts stems while ignoring time.

\begin{enumerate}
\setcounter{enumi}{6}

% Source map: 2025 RC, Question 07 — shared purpose
\item A principal purpose shared by both passages is to
\begin{enumerate}
\item[(A)] argue that cities should stop planting trees.
\item[(B)] explain why urban tree policy must consider more than the number of individual trees.
\item[(C)] show that old trees are always dangerous.
\item[(D)] establish that cultural value is measurable in dollars.
\item[(E)] defend planting a single species throughout a city.
\end{enumerate}

% Source map: 2025 RC, Question 08 — detail unique to passage B
\item Which consideration is mentioned in Passage B but not Passage A?
\begin{enumerate}
\item[(A)] Vulnerability to pests
\item[(B)] Unequal distribution of shade
\item[(C)] Animal shelter created by hollows that form slowly
\item[(D)] Species diversity
\item[(E)] The cost of maintaining a famous tree
\end{enumerate}

% Source map: 2025 RC, Question 09 — detail unique to passage A
\item Which consideration is discussed in Passage A but not Passage B?
\begin{enumerate}
\item[(A)] Residents' stories about landmarks
\item[(B)] The possibility that one pest could devastate closely related trees
\item[(C)] Slow formation of wildlife habitat
\item[(D)] Removal only for a specific hazard
\item[(E)] The hydrological role of extensive roots
\end{enumerate}

% Source map: 2025 RC, Question 10 — comparison of methods
\item The passages develop their positions in similar ways because both
\begin{enumerate}
\item[(A)] acknowledge a value emphasized by the other position before limiting its priority.
\item[(B)] rely chiefly on numerical studies.
\item[(C)] describe the history of one city's planting program.
\item[(D)] argue that two goals can never be pursued together.
\item[(E)] define veteran trees solely by their height.
\end{enumerate}

% Source map: 2025 RC, Question 11 — likely agreement
\item The authors of both passages would most likely agree that
\begin{enumerate}
\item[(A)] every veteran tree must be preserved regardless of danger.
\item[(B)] planting should increase species diversity in urban forests.
\item[(C)] young trees provide no present benefit.
\item[(D)] canopy should be concentrated near civic buildings.
\item[(E)] ecological value and fair distribution are unrelated.
\end{enumerate}

% Source map: 2025 RC, Question 12 — principle difference
\item Which principle is most consistent with Passage A but not Passage B?
\begin{enumerate}
\item[(A)] When resources are limited, system-wide resilience and equitable coverage may justify declining to preserve a valued individual component.
\item[(B)] An old living system can perform functions that replacements will not immediately provide.
\item[(C)] Diversity can be increased while selected veteran trees remain protected.
\item[(D)] Removal should require a specific hazard or unavoidable conflict.
\item[(E)] Cultural attachment can be relevant to public decisions.
\end{enumerate}

% Source map: 2025 RC, Question 13 — cross-passage viewpoint
\item The author of Passage A would most likely respond to Passage B's claim about hollows by saying that
\begin{enumerate}
\item[(A)] hollows have no ecological value.
\item[(B)] every hollow tree is unsafe.
\item[(C)] this value should be considered, but it cannot automatically override canopy-wide resilience and distribution.
\item[(D)] young trees develop hollows immediately.
\item[(E)] cultural stories are more important than wildlife.
\end{enumerate}

% Source map: 2025 RC, Question 14 — application
\item Which policy would the author of Passage B be most likely to support?
\begin{enumerate}
\item[(A)] Replace all trees over fifty years old with one fast-growing species.
\item[(B)] Protect a healthy veteran oak, plant diverse younger trees nearby, and expand planting in low-canopy neighborhoods.
\item[(C)] Stop all planting until every veteran tree has died naturally.
\item[(D)] Preserve veteran trees only when maintenance costs nothing.
\item[(E)] Rank trees exclusively by trunk diameter.
\end{enumerate}

\end{enumerate}

\subsection*{Packet 3: Describing a Community Archive}

The Harbor City Archive is digitizing thousands of photographs contributed by neighborhood organizations. Its catalog currently uses a controlled vocabulary developed decades ago by professional archivists. Consistent terms make searching easier, but some of those terms are unfamiliar to the people represented in the photographs; others flatten distinctions that residents consider important. A single official label, for example, may combine several locally distinct festivals.

The archive's director proposes allowing community organizations to add preferred descriptions alongside the controlled terms. Critics worry that multiple names will fracture the catalog: a researcher who knows one name might miss records filed under another. That concern would be decisive if the proposal replaced controlled vocabulary. It does not. Each record would retain standardized terms while also storing community descriptions as linked fields. A search for either could retrieve the same record, and the links could reveal that two expressions overlap without asserting that they are identical.

The proposal also changes who can contest a description. Under the present system, residents may send corrections, but staff decide whether a controlled term is technically mistaken. A description can therefore remain officially ``correct'' even when it misidentifies how a community understands an event. Under the linked-field proposal, staff would continue to maintain the stable research vocabulary, while contributing organizations would maintain their own descriptive fields. Disagreement would remain visible rather than being resolved by giving one party exclusive authority.

This approach is sometimes criticized as abandoning archival neutrality. Yet the current vocabulary is not neutral simply because its choices are old or professionally standardized. Classification always highlights some similarities and suppresses others. The responsible response is not to abandon standards, which make a large collection usable, but to expose important choices and permit additional well-sourced descriptions. Harbor City should adopt the proposal as a practical form of shared stewardship.

The director should, however, establish an appeal process for harmful or unsupported descriptions. Shared authority does not mean that every submitted phrase must be displayed. The archive can require organizations to identify the source of a preferred name, document disputed changes, and distinguish a group's self-description from a claim about others. These safeguards would make the system accountable without restoring the fiction that only one legitimate vocabulary exists.

\begin{enumerate}
\setcounter{enumi}{14}

% Source map: 2025 RC, Question 15 — main point
\item Which statement best expresses the main point of the passage?
\begin{enumerate}
\item[(A)] Harbor City should discard controlled vocabulary and let each contributor catalog photographs independently.
\item[(B)] Digitization is inappropriate for community photographs.
\item[(C)] Harbor City should retain standardized search terms while adding accountable community-controlled descriptions as linked fields.
\item[(D)] Archival descriptions are neutral whenever trained professionals write them.
\item[(E)] Every proposed description should be displayed without review.
\end{enumerate}

% Source map: 2025 RC, Question 16 — function of detail
\item Why does the author mention that a search for either kind of term could retrieve the same record?
\begin{enumerate}
\item[(A)] To show how linked fields answer the concern that multiple names would make records harder to find
\item[(B)] To prove that all names mean exactly the same thing
\item[(C)] To suggest that controlled terms should be deleted
\item[(D)] To show that residents rarely search the archive
\item[(E)] To explain why digitization is too expensive
\end{enumerate}

% Source map: 2025 RC, Question 17 — inference
\item Which statement can most reasonably be inferred from the passage?
\begin{enumerate}
\item[(A)] The author regards consistent search as valuable even while criticizing exclusive control over description.
\item[(B)] The director believes professional archivists should have no role in the catalog.
\item[(C)] Community organizations always agree with one another.
\item[(D)] The archive currently displays every correction it receives.
\item[(E)] Linked fields necessarily reduce the number of records returned by a search.
\end{enumerate}

% Source map: 2025 RC, Question 18 — phrase function
\item The phrase ``officially correct'' primarily emphasizes that
\begin{enumerate}
\item[(A)] every controlled term is factually false.
\item[(B)] technical conformity to an archival standard may not capture the represented community's understanding.
\item[(C)] residents are forbidden to contact archive staff.
\item[(D)] community terminology never changes.
\item[(E)] only courts can determine correct language.
\end{enumerate}

% Source map: 2025 RC, Question 19 — author agreement
\item The author would most likely agree that
\begin{enumerate}
\item[(A)] visible disagreement in a catalog is always a sign of failure.
\item[(B)] an old classification is neutral unless someone proves malicious intent.
\item[(C)] search consistency and community self-description can be supported by different, linked fields.
\item[(D)] shared stewardship requires publishing unsupported claims.
\item[(E)] professional standards have no practical value.
\end{enumerate}

% Source map: 2025 RC, Question 20 — strengthen
\item Which new information would most strengthen the author's recommendation?
\begin{enumerate}
\item[(A)] A pilot using linked fields let researchers find records through either vocabulary without reducing retrieval accuracy.
\item[(B)] Some photographs are black and white.
\item[(C)] The archive building needs roof repairs.
\item[(D)] One contributor prefers printed catalogs.
\item[(E)] Professional archivists developed the existing vocabulary.
\end{enumerate}

% Source map: 2025 RC, Question 21 — underlying principle
\item Which principle most closely underlies the author's argument?
\begin{enumerate}
\item[(A)] A useful information system may preserve common standards while making room for accountable perspectives those standards omit.
\item[(B)] The oldest description of an object is always the most accurate.
\item[(C)] A catalog should contain only terms understood by every possible user.
\item[(D)] Administrative efficiency should override every other value.
\item[(E)] Giving more than one group authority makes accountability impossible.
\end{enumerate}

\end{enumerate}

\subsection*{Packet 4: The Value of Waiting Seeds}

Managers restoring a damaged grassland often judge a seed mix by how many seeds germinate during the first favorable season. This measure is convenient, but it can reward the wrong property. A population in which nearly every seed sprouts after the same rainfall may produce an impressive first year. If the next drought arrives before those plants reproduce, however, the population has no reserve. By contrast, a species whose seeds germinate across several years may look unsuccessful at first while retaining the capacity to recover.

Delayed germination is not simply inactivity. A seed's response depends on interacting cues such as temperature shifts, moisture duration, light exposure, and changes to its coat. Because the cues combine, two apparently similar storms need not produce the same result. Researchers who test only one cue at a time can therefore miss how seeds behave in the field. Likewise, reporting an average germination time can hide a strategically important spread: an average of two years might describe a population whose seeds all sprout near year two, or one split between immediate sprouting and much longer delay.

This variation can function as a biological form of risk distribution. Seeds that sprout early can exploit a good season; seeds that remain dormant protect the population if that apparent opportunity disappears. The balance is not universally beneficial. In a highly stable environment, prolonged dormancy may merely postpone reproduction, and seeds can die while waiting. But where conditions fluctuate unpredictably, a range of responses can allow a lineage to persist even when any one year's cohort fails.

Restoration practice should therefore measure more than first-season germination. Small samples can be left under realistic field conditions for several years, while laboratory trials test combinations of cues rather than isolated triggers. Managers should also distinguish viable dormant seeds from dead ones. None of these measures guarantees long-term success, since competition, grazing, and changing climate also matter. They do, however, prevent a rapid but fragile response from being mistaken for resilience.

The larger lesson concerns the timescale of evidence. A short observation can accurately describe what happened first without revealing the strategy that produced it. When a system stores some of its possibilities for later, waiting is part of the outcome rather than an absence of outcome.

\begin{enumerate}
\setcounter{enumi}{21}

% Source map: 2025 RC, Question 22 — primary purpose
\item The author is primarily concerned with
\begin{enumerate}
\item[(A)] proving that dormant seeds never die.
\item[(B)] arguing that first-season germination alone can misrepresent resilience and proposing broader assessment.
\item[(C)] showing that laboratory studies should replace field studies.
\item[(D)] identifying one universal germination schedule for all grasslands.
\item[(E)] explaining why rainfall is irrelevant to seeds.
\end{enumerate}

% Source map: 2025 RC, Question 23 — supported statement
\item The passage most strongly supports which statement?
\begin{enumerate}
\item[(A)] Two populations with the same average germination time can differ meaningfully in the spread of their responses.
\item[(B)] Every seed responds to moisture in the same way.
\item[(C)] Immediate germination is always maladaptive.
\item[(D)] A stable environment favors the longest possible dormancy.
\item[(E)] Field observation can identify no useful information.
\end{enumerate}

% Source map: 2025 RC, Question 24 — analogy
\item The author's criticism of relying only on first-season germination is most analogous to criticizing
\begin{enumerate}
\item[(A)] a music teacher for evaluating a rehearsal before learning the musicians' names.
\item[(B)] a librarian for arranging books by subject rather than by height.
\item[(C)] an investor for judging a reserve fund solely by its first month's spending, without considering resources deliberately held for later shocks.
\item[(D)] a chef for tasting soup before serving it.
\item[(E)] a cyclist for checking tire pressure before a trip.
\end{enumerate}

% Source map: 2025 RC, Question 25 — explain characterization
\item Why does the author call first-season germination a measure that can ``reward the wrong property''?
\begin{enumerate}
\item[(A)] It may favor rapid, synchronized sprouting even when that leaves no reserve for later failure.
\item[(B)] It counts only seeds that are dead.
\item[(C)] It requires managers to wait too many years.
\item[(D)] It measures competition rather than germination.
\item[(E)] It assumes that grasslands never experience rain.
\end{enumerate}

% Source map: 2025 RC, Question 26 — example function
\item The comparison between two populations with a two-year average primarily serves to
\begin{enumerate}
\item[(A)] establish that averages are always mathematically incorrect.
\item[(B)] illustrate how an average can conceal variation relevant to a population's strategy.
\item[(C)] prove that all seeds should wait exactly two years.
\item[(D)] show that laboratory trials are unnecessary.
\item[(E)] identify the maximum lifespan of a seed.
\end{enumerate}

% Source map: 2025 RC, Question 27 — inference
\item Which statement can most reasonably be inferred from the passage?
\begin{enumerate}
\item[(A)] A restoration with lower initial germination may sometimes be more resilient than one with higher initial germination.
\item[(B)] Dormancy guarantees survival under every climate.
\item[(C)] Managers can ignore competition and grazing once dormancy is measured.
\item[(D)] Seeds that germinate early never reproduce.
\item[(E)] Only dead seeds fail to germinate in the first year.
\end{enumerate}

\end{enumerate}

\subsection*{Answer Key and Concise Explanations}

\subsubsection*{Packet 1}
\begin{enumerate}
\item \textbf{B.} The passage argues for making acoustic access visible in decisions without pretending it has one objective score.
\item \textbf{D.} The passage assigns equipment measurement, not instrument calibration, to resident annotations.
\item \textbf{B.} Vale accepts variation while defending narrower review questions grounded in evidence.
\item \textbf{A.} The qualification limits the claim but leaves the policy's earlier, institutional benefit intact.
\item \textbf{D.} It combines early review, access cues, and evidence about actual use.
\item \textbf{A.} The passage treats the two evidence types as complementary and individually incomplete.
\end{enumerate}

\subsubsection*{Packet 2}
\begin{enumerate}
\setcounter{enumi}{6}
\item \textbf{B.} Each author argues that counting trees alone omits system-level or time-dependent values.
\item \textbf{C.} Only Passage B discusses slowly formed hollows as animal habitat.
\item \textbf{B.} Only Passage A explicitly raises common vulnerability to a single pest.
\item \textbf{A.} A concedes old-tree value; B concedes the need for diversity and broader planting.
\item \textbf{B.} Both expressly endorse greater species diversity.
\item \textbf{A.} Passage A permits this tradeoff; Passage B resists treating valued veterans as expendable system units.
\item \textbf{C.} This response acknowledges the function while preserving Passage A's priority structure.
\item \textbf{B.} It combines veteran-tree preservation, diversity, and fair geographic expansion.
\end{enumerate}

\subsubsection*{Packet 3}
\begin{enumerate}
\setcounter{enumi}{14}
\item \textbf{C.} The proposal preserves standards while adding sourced, community-maintained linked descriptions.
\item \textbf{A.} Cross-retrieval directly addresses discoverability under multiple names.
\item \textbf{A.} The recommended system preserves the stable vocabulary precisely because consistency remains useful.
\item \textbf{B.} The quotation marks distinguish technical correctness from adequate community identification.
\item \textbf{C.} Linked fields let the two goals coexist without pretending the vocabularies are identical.
\item \textbf{A.} Successful pilot evidence supports both usability and inclusive description.
\item \textbf{A.} This principle captures the passage's standards-plus-accountability position.
\end{enumerate}

\subsubsection*{Packet 4}
\begin{enumerate}
\setcounter{enumi}{21}
\item \textbf{B.} The passage critiques a short-term metric and recommends multi-year, multi-cue evidence.
\item \textbf{A.} The passage explicitly contrasts equal averages with different response distributions.
\item \textbf{C.} In both cases, resources held for later are part of resilience but disappear from a short initial measure.
\item \textbf{A.} Synchronized early sprouting can create a strong first year and a fragile population.
\item \textbf{B.} The example shows that a shared average need not represent a shared pattern.
\item \textbf{A.} The argument repeatedly allows a slower initial response to preserve later recovery capacity.
\end{enumerate}
